%Master File:lectures.tex
\newcommand{\Kb}{\bar{K}}

\lesson{Population Genetics}
\vspace{-1cm}
\begin{center}
  \includegraphics[height=10cm]{\figs/GAroulette/study.png}
\end{center}
\keywords{Fisher-Wright model, ODEs, Markov chains.}
%%%%%%%%%%%%%%%%%%%%%%% Next Slide %%%%%%%%%%%%%%%%%%%%%%%
\renewcommand{\Outline}{%
\begin{slide}
\section[1]{Outline}

\begin{minipage}{11cm}
  \begin{enumerate}
    \outlineitem{Population Genetics}{populations}
    \outlineitem{Facts of Life}{fisherwright}
    \outlineitem{Infinite Population Model}{infpop}
    \outlineitem{Finite Population Model}{finpop}
  \end{enumerate}
\end{minipage}\hfill
\begin{minipage}{11.5cm}
  \includegraphics[width=\linewidth]{\figs/GAroulette/study.png}
\end{minipage}
\end{slide}
\addtocounter{outlineitem}{1}
}

\setcounter{outlineitem}{1}
\Outline % Modelling
\toptarget{firstoutline}
%%%%%%%%%%%%%%%%%%%%%%% Next Slide %%%%%%%%%%%%%%%%%%%%%%%

\begin{slide}
\section[-1]{Evolution and Molecular Biology}

\begin{PauseHighLight}
  \begin{itemize}
  \item Evolution is the theory that underpins our understanding of
    biology\pause
  \item Molecular biology has open up a whole new perspective on
    evolution as we see similar proteins in related species\pause
  \item This evolution ancestry helps us understand molecular biology
    as similar proteins usually share the same function\pause
  \item However we can also use molecular biology to help understand
    evolution---more of this when we discuss phylogenetic tree\pause
  \item But how does evolution work?\pause
  \end{itemize}
\end{PauseHighLight}

\end{slide}

%%%%%%%%%%%%%%%%%%%%%%% Next Slide %%%%%%%%%%%%%%%%%%%%%%%

\begin{slide}
\section{Evolutionary Population Dynamics}

\begin{PauseHighLight}
  \begin{itemize}
  \item Since as least the 1920's researchers have tried to model
    evolution\pause
  \item The field is known as \emph{population genetics}\pause
  \item Much of it was developed before we understood the mechanisms of
    heredity\pause
  \item It is therefore rather abstracted from biology\pause
  \item However it was enormously influential to our understanding of
    evolution\pause
  \end{itemize}
\end{PauseHighLight}

\end{slide}


%%%%%%%%%%%%%%%%%%%%%%% Next Slide %%%%%%%%%%%%%%%%%%%%%%%
\Outline % Fisher-Wright Model
%%%%%%%%%%%%%%%%%%%%%%% Next Slide %%%%%%%%%%%%%%%%%%%%%%%

\begin{slide}
\section[-2]{Fisher-Wright Model}

\begin{PauseHighLight}
  \begin{itemize}
  \item The classic model of population genetics is the Fisher-Wright
    (or Wright-Fisher) model developed independently by Fisher and
    Wright around 1930\pause
  \item It is important because it laid the theoretical framework for
    many of the prejudices held today\pause
  \item In its simplest form it is a model of what happens at a single
    locus (site) under selection, mutation and crossover
  \end{itemize}
  \begin{center}
    \includegraphics[height=4cm]{loci}\pause
  \end{center}
\end{PauseHighLight}

\end{slide}

%%%%%%%%%%%%%%%%%%%%%%% Next Slide %%%%%%%%%%%%%%%%%%%%%%%

\begin{slide}
\section{Locus and Alleles}

\begin{PauseHighLight}
  \begin{itemize}
  \item The locus can be thought of as the location of a gene\pause
  \item The type of gene at a site will differ between individuals
    (e.g. one for smooth peas and one for wrinkled peas)\pause
  \item The different types of genes are known as \emph{alleles}\pause
  \item We assume that genes can be modified from one allele to another
    by mutation\pause
  \item Alleles are also exchanged during sexual reproduction
    (crossover)\pause 
  \end{itemize}
\end{PauseHighLight}

\end{slide}

%%%%%%%%%%%%%%%%%%%%%%% Next Slide %%%%%%%%%%%%%%%%%%%%%%%

\begin{slide}
\section{Fitness}

\begin{PauseHighLight}
  \begin{itemize}
  \item Fitness has a technical meaning in population genetics; it
    is a measure of the expected number of offspring an individual
    produces\pause
  \item If one individual is 10\% fitter than another it will (in
    expectation) produce 10\% more offspring\pause
  \item Of course, the individual may get flattened by a car before it
    reaches sexual maturity and never produce any offspring\pause
  \item Thus fitness measures potential to reproduce\pause
  \end{itemize}
\end{PauseHighLight}

\end{slide}


%%%%%%%%%%%%%%%%%%%%%%% Next Slide %%%%%%%%%%%%%%%%%%%%%%%

\begin{slide}
\section{Epistatic Interaction}

\begin{PauseHighLight}
  \begin{itemize}
  \item Many genes interact (we now know they often code for proteins
    that directly interact chemically in many cases)\pause
  \item Thus the fitness of an allele will typically depend on other
    alleles at other locations\pause
  \item This interaction is known as \emph{epistasis}\pause
  \item To keep things simple in the Fisher-Wright model we assume no
    epistasis\pause
  \item There exists more complex models which treat epistasis between
    two or more loci, but these rapidly become rather complex\pause    
  \end{itemize}
\end{PauseHighLight}

\end{slide}

%%%%%%%%%%%%%%%%%%%%%%% Next Slide %%%%%%%%%%%%%%%%%%%%%%%

\begin{slide}
\section[-1]{Sex}

\begin{PauseHighLight}
  \begin{itemize}
  \item Sex is surprisingly ubiquitous in higher organisms\pause
  \item Few higher organisms seem to reproduce asexually (although
    even if they do, many sometimes also reproduce sexually)\pause
  \item Asexual higher organisms all would have produced sexually
    sometime in their past\pause
  \item Why sex is surprising is that if females produce asexually
    (producing clones) then they breed at twice the rate of mixed
    couples\pause---hence they are twice as fit\pauseb
  \item Sexual reproduction has to give at least a two fold advantage in
    fitness on a very short time scale to survive\pause
  \end{itemize}
\end{PauseHighLight}

\end{slide}

%%%%%%%%%%%%%%%%%%%%%%% Next Slide %%%%%%%%%%%%%%%%%%%%%%%

\begin{slide}
\section[-1]{Linkage Equilibrium}

\begin{PauseHighLight}
  \begin{itemize}\squeeze
  \item As well as no epistasis the model assumes lots of crossover
    (compared to selection and mutation)\pause
  \item This is not unrealistic in many cases\pause
  \item Under this assumption we can treat the loci independently and
    concentrate on a single locus\pause
  \item This is known as \emph{linkage equilibrium}\pause
  \item We therefore only need to consider the proportion of the
    different alleles at the locus\pause
  \item In an asexual population, the alleles at different sites are
    coupled and the situation is much harder (in fact, unsolved)\pause
  \end{itemize}
\end{PauseHighLight}

\end{slide}

%%%%%%%%%%%%%%%%%%%%%%% Next Slide %%%%%%%%%%%%%%%%%%%%%%%

\begin{slide}
\section{Diploidy}

\begin{PauseHighLight}
  \begin{itemize}
  \item Most organisms are diploids, that is they carry two copies of
    DNA; one from their father and one from their mother\pause
  \item Like sex, it is not altogether clear what advantage diploidy
    brings\pause
  \item Clearly, diploids can tolerate mutations that make one set of
    genes dysfunctional\pause
  \item But, this also allows detrimental mutations to build up in a
    genome\pause
  \item Incidentally, mating between close relatives is likely to create
    nonviable offspring because it is likely to combine two dysfunctional
    alleles\pause
  \end{itemize}
\end{PauseHighLight}

\end{slide}


%%%%%%%%%%%%%%%%%%%%%%% Next Slide %%%%%%%%%%%%%%%%%%%%%%%

\begin{slide}
\section[-1]{Haploidy}

\begin{PauseHighLight}
  \begin{itemize}
  \item To keep things simple we are going to assume each individual
    only carries one set of genes\pause
  \item Organisms that do this are said to be \emph{haploidy}\pause
  \item There are various organisms (e.g. male bees, wasps, ants, many
    fungi) that are haploid\pause
  \item Generalising to diploid organisms is just fiddly\pause---it
    mainly involves adding in factors of two every so often\pauseb,
    but it does not substantially change things\pauseb
  \item Incidentally many plants are polypoids having more than two
    copies of each gene\pause
  \end{itemize}
\end{PauseHighLight}

\end{slide}

%%%%%%%%%%%%%%%%%%%%%%% Next Slide %%%%%%%%%%%%%%%%%%%%%%%

\begin{slide}
\section{Gender}

\begin{PauseHighLight}
  \begin{itemize}
  \item Life tends also to be gendered\pause
  \item This possibly just reflects that finding a mate consumes energy
    and by specialisation some energy can be saved\pause
  \item However, this is another complicating factor, which adds little
    to the model\pause
  \item I will assume a \emph{monoecious} population, i.e. one with only
    one gender, although it still mates\pause
  \item These are all simplifying assumptions, but it allows us to
    concentrate on the role of selection, mutation and crossover\pause
  \end{itemize}
\end{PauseHighLight}

\end{slide}



%%%%%%%%%%%%%%%%%%%%%%% Next Slide %%%%%%%%%%%%%%%%%%%%%%%
\Outline % Infinite Size Population
%%%%%%%%%%%%%%%%%%%%%%% Next Slide %%%%%%%%%%%%%%%%%%%%%%%

\begin{slide}
\section[-1]{Single Locus Model}

\begin{PauseHighLight}
  \begin{itemize}
  \item We consider the case where the allele can take just two states
    $A$ and $B$\pause
  \item We call $B$ the mutant\pause
  \item We make the following assumptions
    \begin{itemize}\squeeze
    \item $B$ has a selective advantage of $1+s$\pause
    \item $A$ mutates to $B$ with a forward mutation rate of $u$\pause
    \item $B$ mutates to $A$ with a backward mutation rate of $v$\pause
    \item The population size, $P$, is unchanged from generation to generation\pause
    \end{itemize}
  \end{itemize}
\end{PauseHighLight}

\end{slide}

%%%%%%%%%%%%%%%%%%%%%%% Next Slide %%%%%%%%%%%%%%%%%%%%%%%

\begin{slide}
\section[-1]{Many Seeds}

\begin{PauseHighLight}
  \begin{itemize}
  \item We consider the population changing from one generation to the
    next by a three stage process\pause
    \begin{itemize}
    \item Each individual produces a very large number of seeds in
      proportional to its fitness (selective advantage)\pause
    \item The seeds are mutated according to the mutation
      probabilities\pause
    \item We select a sample of $P$ of the seeds independently to form
      the next generation\pause
    \end{itemize}
  \item Initially we assume $P$ is infinite and consider only the
    proportion of mutants in the population\pause
  \end{itemize}
\end{PauseHighLight}
\end{slide}

%%%%%%%%%%%%%%%%%%%%%%% Next Slide %%%%%%%%%%%%%%%%%%%%%%%

\begin{slide}
\section[-0.5]{Dynamics}

\begin{PauseHighLight}
  \begin{itemize}
  \item Assume we start with $n$ mutants\pause
  \item After selection we have a proportion $p_s$ of mutant seeds
    \begin{align*}
      p_s = \frac{(1+s)n}{(1+s)n + P-n}\pause = \frac{(1+s)n}{P+s\,n}\pause
    \end{align*}
  \item After mutation we have a proportion $p_{sm}$ of mutants
    \begin{align*}
      p_{sm} = (1-v) p_s + u (1-p_s)\pause = \frac{(1-v)(1+s)n +
        u(P-n)}{P+s\,n}\pause
    \end{align*}
  \end{itemize}
\end{PauseHighLight}

\end{slide}

%%%%%%%%%%%%%%%%%%%%%%% Next Slide %%%%%%%%%%%%%%%%%%%%%%%

\begin{slide}
\section{Infinite Population Approximation}

\begin{PauseHighLight}
  \begin{itemize}
  \item If we assume an infinite population then this will be the
    proportion of individuals at the next generation
    \begin{align*}
      p(t+1) = \frac{(1-v)(1+s) p(t) +
        u\bigl(1-p(t)\bigr)}{1+s\,p(t)}\pause
    \end{align*}
  \item If the effect of selection and mutation is small then we can
    approximate the dynamics by a differential equation
    \begin{align*}
      \frac{\dd p}{\dd t} \approx p(t+1) - p(t) = 
      \frac{u - (u+v-s+v\,s) p(t) - s\, p^2(t)}{1+s\,p(t)}\pause
    \end{align*}
  \end{itemize}
\end{PauseHighLight}

\end{slide}

%%%%%%%%%%%%%%%%%%%%%%% Next Slide %%%%%%%%%%%%%%%%%%%%%%%

\begin{slide}
\section[-1]{Infinite Population Prediction}

\begin{center}
  \includegraphics[height=14cm]{infpop}\\
  $s=0.1$, $u=v=0.001$
\end{center}
\end{slide}

%%%%%%%%%%%%%%%%%%%%%%% Next Slide %%%%%%%%%%%%%%%%%%%%%%%

\begin{slide}
\section{Difficulties of Infinite Populations}

\begin{PauseHighLight}
  \begin{itemize}
  \item In an infinite population a mutation (even if it is very
    unlikely) will always occur in a small fraction of individuals\pause
  \item Selection can rapidly amplify the mutants if they are
    fitter---note fitter individuals initial grow exponentially
    \begin{align*}
      (1+s)^t = \e{t\,\log(1+s)} \approx \e{s\,t}\pause
    \end{align*}
  \item If the mutation rate is small then in finite populations the
    mutant may take a very long time to appear\pause
  \item Thus infinite populations can have far more rapid take over
    times than is experience in large, but finite, populations\pause
  \end{itemize}
\end{PauseHighLight}

\end{slide}


%%%%%%%%%%%%%%%%%%%%%%% Next Slide %%%%%%%%%%%%%%%%%%%%%%%
\Outline % Finite Population Models
%%%%%%%%%%%%%%%%%%%%%%% Next Slide %%%%%%%%%%%%%%%%%%%%%%%

\begin{slide}
\section[-2]{Finite Populations}

\begin{PauseHighLight}
  \begin{itemize}
  \item To model finite populations, we have to draw a sample of $P$
    individuals from the seeds\pause
  \item The probability of drawing $n'=n(t+1)$ mutants given a proportion
    $p_{sm}(n(t))$ of mutants in the seed population is given by a binomial
    distribution
    \begin{align*}
      \Prob{n'} = \binom{P}{n'} 
      p_{sm}^{n'}
      \left(\strut 1 - p_{sm}\right)^{P-n'}\pause
    \end{align*}
  \item (Recall that $p_{sm}$ is the proportion of mutant seeds
    after selection and mutation and depends on $n(t)$)\pause
  \item We can simulate the evolution of a population, by drawing a
    random binomial deviate at each generation\pause
  \end{itemize}
\end{PauseHighLight}

\end{slide}

%%%%%%%%%%%%%%%%%%%%%%% Next Slide %%%%%%%%%%%%%%%%%%%%%%%

\begin{slide}
\section{Finite Population Run}

\begin{center}
  \includegraphics[height=14cm]{finitepop}\\
  $s=0.1$, $u=v=0.001$
\end{center}
\end{slide}


%%%%%%%%%%%%%%%%%%%%%%% Next Slide %%%%%%%%%%%%%%%%%%%%%%%

\begin{slide}
\section[-2]{Markov Models}
\toptarget{markov}

\begin{PauseHighLight}
  \begin{itemize}\squeeze
  \item The evolving population can be modelled as a Markov chain\pause
  \item The population at time $t$ depends only on the population at
    time $t-1$\pause
  \item It is \textit{conditionally independent} of the population at
    earlier times \textit{given} the population at time $t-1$\pause
  \item We define $p_n(t)=\Prob{n(t)=n}$ to be the probability that
    there are $n$ mutants at time $t$\pause
  \item Let $W(n\rightarrow n')$ be the transition probability of a
    population with $n$ mutants going to a population of $n'$ mutants in
    the next generation\pause
    \begin{align*}
      p_{n'}(t+1) = \sum_{n=0}^P W(n\rightarrow n') p_n(t)\pause
    \end{align*}
  \end{itemize}
\end{PauseHighLight}

\end{slide}

%%%%%%%%%%%%%%%%%%%%%%% Next Slide %%%%%%%%%%%%%%%%%%%%%%%

\begin{slide}
\section[-1]{Matrix Formulation}

\begin{PauseHighLight}
  \begin{itemize}
  \item Denote the \textit{probability vector} $(p_0(t), p_1(t), \ldots,
    p_P(t))$ by $\bm{p}(t)$\pause
  \item Let $\mat{W}$ be the \textit{transition matrix} with elements
    $W_{n' n} =W(n\rightarrow n')$\pause
  \item The dynamics can be written in matrix form
    \begin{align*}
      \bm{p}(t) = \mat{W} \bm{p}(t-1) \pause
      = \mat{W} \left( \strut \mat{W} \bm{p}(t-2) \right) \pauseb
      = \mat{W} \mat{W} \cdots \mat{W} \bm{p}(0)\pauseb
      = \mat{W}^t \bm{p}(0)\pauseb
    \end{align*}
  \item Where $\bm{p}(t=0)$ encodes our \textit{initial conditions}\pause
  \item If initially $n(t=0)=0$ then $p_0(t=0)=1$ and $p_i(t=0)=0, \forall
    i>0$\pause
  \end{itemize}
\end{PauseHighLight}

\end{slide}



%%%%%%%%%%%%%%%%%%%%%%% Next Slide %%%%%%%%%%%%%%%%%%%%%%%

\begin{slide}
\section[-1]{Markov Analysis}

\begin{PauseHighLight}
  \begin{itemize}
  \item In our problem the transition matrix, $\mat{W}$, has elements
    \begin{align*}
      W_{n' n} = \binom{P}{n'}\left(\strut p_{sm}(n)\right)^{n'}
      \left(\strut 1 - p_{sm}(n) \right)^{P-n'}\pause
    \end{align*}
  \item where $p_{sm}(n)$ is given by
    \begin{align*}
      p_{sm}(n) = \frac{(1-v)(1+s)n + u(P-n)}{P+s\,n}\pause
    \end{align*}
  \item This provides an exact formulation of a finite population\pause
  \end{itemize}
\end{PauseHighLight}

\end{slide}

%%%%%%%%%%%%%%%%%%%%%%% Next Slide %%%%%%%%%%%%%%%%%%%%%%%

\begin{slide}
\section[-2]{Example of Markov Analysis}
\pause
\pb
\pauselevel{=1}
\begin{center}
  \multipdf[width=\linewidth]{MarkovFW}\pause
\end{center}

\end{slide}

%%%%%%%%%%%%%%%%%%%%%%% Next Slide %%%%%%%%%%%%%%%%%%%%%%%

\begin{slide}
\section{Mean Number of Mutants}
\begin{center}
  \includegraphics[width=\linewidth]{markov_full}
\end{center}
\end{slide}

%%%%%%%%%%%%%%%%%%%%%%% Next Slide %%%%%%%%%%%%%%%%%%%%%%%

\begin{slide}
\section{Problems with Markov Chain Analysis}

\begin{PauseHighLight}
  \begin{itemize}
  \item The Markov Chain analysis is exact and provides a full picture
    of the probability of the population being in a certain state\pause
  \item Unfortunately, the transition matrix $W(n\rightarrow n')$
    is of size the number of states squared\pause
  \item The number of states for a population of size $P$ is $P+1$\pause
  \item For a large population writing the transition matrix becomes
    infeasible\pause
  \item Even when it is feasible we don't get much intuition on changing
    the parameters\pause
  \end{itemize}
\end{PauseHighLight}

\end{slide}


%%%%%%%%%%%%%%%%%%%%%%% Next Slide %%%%%%%%%%%%%%%%%%%%%%%

\begin{slide}
\section[-2]{Conclusion}

\begin{PauseHighLight}
  \begin{itemize}
  \item The Fisher-Wright model provides a simple model undergoing
    selection, mutation and crossover\pause
  \item It allows us to explore the effects of selection pressure and
    mutation rate\pause
  \item We can solve the infinite population model directly or by
    solving an ordinary differential equation (ODE)\pause
  \item Unfortunately infinite population models don't always capture
    the behaviour of finite populations even when the population is very
    large\pause
  \item We can simulate finite population models\pause
  \item For small populations we can use Markov models\pause
  \end{itemize}
\end{PauseHighLight}

\end{slide}

